\section*{الفصل \ref{ch:g10:storms-and-cyclones} --- العواصف والأعاصير المدارية}

\begin{solution}{exo:g10:storms-and-cyclones:1}
\omterm{def:g10:storms-and-cyclones:thunderstorm}{العاصفة الرعدية} عاصفة ينتجها ركام مزني، يصحبها \omterm{def:g10:storms-and-cyclones:lightning}{برق} ورعد، وعادة زخات غزيرة وهبات
ريح، وأحيانًا بَرَد. \omterm{def:g10:storms-and-cyclones:lightning}{والبرق} شرارة كهربائية هائلة تفرّغ المناطق المشحونة من
العاصفة. والرعد صوت هواء قناة \omterm{def:g10:storms-and-cyclones:lightning}{البرق}، الذي يسخن فجأة إلى حد يتمدد معه
بصدمة.
\end{solution}

\begin{solution}{exo:g10:storms-and-cyclones:2}
\omterm{def:g10:storms-and-cyclones:tropical-cyclone}{الإعصار المداري} عاصفة يبلغ عرضها مئات الكيلومترات، تتكون فوق بحر مداري دافئ حول
مركز ضغط منخفض، رياحها تلتف حلزونيًا نحو الداخل وتحيط بها أشرطة من \omterm{def:g10:storms-and-cyclones:thunderstorm}{العواصف
الرعدية}. \omterm{def:g10:storms-and-cyclones:tropical-cyclone}{والعين} مركزه الهادئ الخالي من السحب تقريبًا؛ \omterm{def:g10:storms-and-cyclones:tropical-cyclone}{وجدار العين} حلقة \omterm{def:g10:storms-and-cyclones:thunderstorm}{العواصف
الرعدية} الشاهقة المحيطة \omterm{def:g10:storms-and-cyclones:tropical-cyclone}{بالعين}. وتكون الرياح على أشدها في \omterm{def:g10:storms-and-cyclones:tropical-cyclone}{جدار العين}.
\end{solution}

\begin{solution}{exo:g10:storms-and-cyclones:3}
\num{70}~عقدة: الفئة 1. و\num{120}~عقدة: الفئة 4.
و\num{140}~عقدة: الفئة 5. والأخيران من الأعاصير الكبرى (الفئة 3
فما فوق).
\end{solution}

\begin{solution}{exo:g10:storms-and-cyclones:4}
$343 \times 9 \approx \qty{3.1}{km}$، ثم $343 \times 15 \approx
\qty{5.1}{km}$. فالوميض الثاني أبعد: العاصفة تبتعد.
\end{solution}

\begin{solution}{exo:g10:storms-and-cyclones:5}
يميل \omterm{def:g10:storms-and-cyclones:lightning}{البرق} إلى ضرب أعلى نقطة في المكان: والشجرة القائمة وحدها في حقل هي تلك
النقطة، فيسري تيار الضربة فيها نزولًا وينتشر في الأرض من حولها. ادخل مبنى أو
سيارة مغلقة، وابقَ فيها إلى ما بعد آخر رعد بوقت طويل.
\end{solution}

\begin{solution}{exo:g10:storms-and-cyclones:6}
العقدة $1852/3600 \approx \qty{0.514}{m/s}$. \num{64}~عقدة: $64 \times
0.514 \approx \qty{32.9}{m/s}$، و$32.9 \times 3.6 \approx
\qty{119}{km/h}$. \num{137}~عقدة: $\qty{70.5}{m/s}$، أي نحو
\qty{254}{km/h} (ويعطي الجدول الرسمي، المقرّب عبر وحدة أخرى،
\qty{252}{km/h}).
\end{solution}

\begin{solution}{exo:g10:storms-and-cyclones:7}
$44 \times 86\,400 \approx \num{3.8e6}$ ومضة في اليوم، و$3.8\times10^6
\times 365.25 \approx \num{1.4e9}$، أي نحو 1.4 مليار في السنة.
\end{solution}

\begin{solution}{exo:g10:storms-and-cyclones:8}
$\Delta p = 1010 - 935 = \qty{75}{hPa}$: انتفاخ قدره نحو \qty{75}{cm}.
\end{solution}

\begin{solution}{exo:g10:storms-and-cyclones:9}
تتكون فوق المحيطات المدارية الدافئة، في الغالب بين خطي العرض \qty{10}{\degree}
و\qty{30}{\degree} في نصفي الكرة، وقليل جدًا منها في جنوب المحيط الأطلسي. وتنجرف
غربًا، ثم تنعطف بعيدًا نحو القطبين، حيث تضعف فوق ماء أبرد. وتغيب عن خط الاستواء
لأن دوران الأرض لا يعطي هناك الهواء المتدفق أي انعطاف: ودونه يتدفق الهواء مباشرة
إلى الضغط المنخفض فيملؤه، ولا يمكن لعاصفة دوّارة أن تبدأ.
\end{solution}

\begin{solution}{exo:g10:storms-and-cyclones:10}
$1031 - 1.56 \times 70 \approx \qty{922}{hPa}$؛ $1031 - 1.56 \times 40
\approx \qty{969}{hPa}$.
\end{solution}

\begin{solution}{exo:g10:storms-and-cyclones:11}
فوق اليابسة تفقد العاصفة إمدادها بالماء المتبخر، ومعه \omterm{def:g10:earth-energy-budget:latent-heat}{الحرارة الكامنة} التي
تحركها؛ ويبطئ السطح الخشن الرياح أيضًا. وفوق الماء البارد يكون التبخر أضعف بكثير،
لأن الهواء البارد يحمل عند التشبع بخار ماء أقل بكثير: فيتضاءل الإمداد \omterm{def:g10:earth-energy-budget:latent-heat}{بالحرارة
الكامنة}.
\end{solution}

\begin{solution}{exo:g10:storms-and-cyclones:12}
عند \qty{0}{\degreeCelsius}: $\sqrt{1.4 \times 8.314 \times 273.15/0.02897}
\approx \qty{331}{m/s}$؛ وعند \qty{30}{\degreeCelsius}، مع $T =
\qty{303.15}{K}$: \qty{349}{m/s}. ومن أجل \qty{9}{s}: $9 \times 331 \approx
\qty{2980}{m}$ و$9 \times 349 \approx \qty{3140}{m}$، أي فرق قدره نحو
\qty{160}{m}، قرابة \qty{5}{\percent}. والقاعدة تعطي \qty{3}{km}: وهذا كافٍ لمعرفة
مدى قرب العاصفة، وهو كل ما وُضعت له.
\end{solution}

\begin{solution}{exo:g10:storms-and-cyclones:13}
$\Delta p = 1010 - 940 = \qty{70}{hPa}$: الانتفاخ نحو \qty{0.7}{m}، وحصة
الريح $4.0 - 0.7 = \qty{3.3}{m}$، أي أكثر من أربعة أضعافه. وفي المد العالي يبلغ
الماء $4.0 + 1.5 = \qty{5.5}{m}$ فوق متوسط مستوى البحر.
\end{solution}

\begin{solution}{exo:g10:storms-and-cyclones:14}
\omterm{def:g10:storms-and-cyclones:tornado}{الإعصار القمعي}: الدورة الواحدة $2\pi \times 100 \approx \qty{628}{m}$، تُقطع في
$628/90 \approx \qty{7}{s}$. والإعصار: $2\pi \times 30\,000 \approx
\qty{188}{km}$، تُقطع في $188\,000/60 \approx \qty{3140}{s}$، أي نحو
\qty{52}{min}. \omterm{def:g10:storms-and-cyclones:tornado}{فالإعصار القمعي} يدور أسرع بنحو 450 مرة: إنه دوامة صغيرة سريعة
جدًا، والإعصار دوامة هائلة أبطأ.
\end{solution}

\begin{solution}{exo:g10:storms-and-cyclones:15}
$24/3888 \approx \qty{0.6}{\percent}$؛ $414/3888 \approx
\qty{10.6}{\percent}$؛ $121/3888 \approx \qty{3.1}{\percent}$. ولو لم يكن لخط
العرض أثر، لبدأ نحو \qty{9}{\percent} من العواصف، أي قرابة 350، على بعد
\qty{5}{\degree} من خط الاستواء؛ ولا يبدأ إلا 24، أي أقل بنحو خمس عشرة مرة. فقرب
خط الاستواء يكون دوران الأرض أضعف من أن يطلق الدوران.
\end{solution}

\begin{solution}{pb:g10:storms-and-cyclones:1}
\textbf{1.} $\pi \times (665\,000)^2 \approx \qty{1.39e12}{m^2}$.

\textbf{2.} $0.015 \times 1.39\times10^{12} \approx \qty{2.08e10}{m^3}$.

\textbf{3.} الكيلومتر المكعب \qty{e9}{m^3}: أي نحو \qty{21}{km^3}.

\textbf{4.} $2.08\times10^{10} \times 1000 \approx \qty{2.08e13}{kg}$.

\textbf{5.} $2.08\times10^{13} \times 2.5\times10^6 \approx \qty{5.2e19}{J}$.

\textbf{6.} $5.2\times10^{19}/86\,400 \approx \qty{6.0e14}{W}$.

\textbf{7.} $6.0\times10^{14}/1.22\times10^{17} \approx 0.005$، أي نحو
\qty{0.5}{\percent}.

\textbf{8.} $240 \times 1.39\times10^{12} \approx \qty{3.3e14}{W}$. فالعاصفة
تحرر من الحرارة قرابة ضعف ضوء الشمس الذي تمتصه مساحتها كلها.

\textbf{9.} لا تعمل العاصفة بضوء الشمس الساقط عليها: بل تحرر حرارة مخزونة على شكل
بخار ماء، تبخر من مساحة أوسع من البحر الدافئ وسحبته الرياح نحو الداخل. فالعاصفة
تركّز طاقة جُمعت فوق منطقة أوسع بكثير.

\textbf{10.} الكتلة نفسها التي للمطر، نحو \qty{2.1e13}{kg} في اليوم (إذ يجب أن
يتجدد البخار المجلوب بالسرعة نفسها التي يسقط بها مطرًا).

\textbf{11.} $6.0\times10^{14}/82 \approx \qty{7.3e12}{m^2}$، أي نحو خمسة أضعاف
مساحة العاصفة: فعلى العاصفة أن تسحب البخار من مساحة أوسع بكثير، ورياحها تجعل البحر
تحتها يتبخر أسرع بكثير من المتوسط.

\textbf{12.} \qty{24.1}{g/kg} عند \qty{28}{\degreeCelsius} مقابل
\qty{14.7}{g/kg} عند \qty{20}{\degreeCelsius}: $24.1/14.7 \approx 1.64$، أي أكثر
بنحو \qty{64}{\percent}.

\textbf{13.} يتبخر الماء الدافئ بسرعة ويحمل الهواء الدافئ بخارًا أكثر بكثير: فيحمل
الهواء المتدفق نحو الداخل \omterm{def:g10:earth-energy-budget:latent-heat}{حرارة كامنة} أكثر بكثير ليحررها، وحين يتشبع يصعد بحرية.
وفوق البحار الأبرد يكون الإمداد أصغر من أن يدير المحرك.

\textbf{14.} فوق اليابسة لا بحر يتبخر: فيتوقف الإمداد \omterm{def:g10:earth-energy-budget:latent-heat}{بالحرارة
الكامنة}، وتبطئ الأرض الخشنة الرياح.

\textbf{15.} $1.5\times10^{12} \times 86\,400 \approx \qty{1.3e17}{J}$.

\textbf{16.} $96 \times 1852/3600 \approx \qty{49.4}{m/s}$.

\textbf{17.} $\frac12 \times 1.16 \times 49.4^2 \approx \qty{1.4e3}{J}$.

\textbf{18.} $5.2\times10^{19}/1.3\times10^{17} = 400$: \omterm{def:g10:earth-energy-budget:latent-heat}{فالحرارة الكامنة} تقارب
400 ضعف الطاقة المعطاة للرياح.

\textbf{19.} معظمها يدفئ الهواء العلوي للعاصفة، الذي ينتشر عند القمة ويطرح الحرارة
نحو الفضاء على شكل أشعة تحت حمراء؛ أما الرياح نفسها فتفقد طاقتها بالاحتكاك بالبحر
والأرض.

\textbf{20.} $1/400 = \qty{0.25}{\percent}$ (ويعطي البرنامج نسبة قدرها
402): \textbf{نحو جزء من 400 من الحرارة التي يحررها التكاثف
يصير ريحًا}.
\end{solution}
